% 付録：本編の置き換えを行列で捉え直す

\section{漸化式と行列}\label{節：特性方程式を行列で見直す}

本編で作った\bookterm{geometric}{等比数列}を,\bookterm{matrix}{行列}で見直す.
\sectionref{節：特性方程式型}の\bookterm{fixed}{不動点}から,一次分数型,隣接3項間型,連立型へ進もう.
以下の「二階」は,次の項に直前の二項が必要なことをいい,本編の隣接3項間型に対応する.

\subsection{一つの値を更新する}\label{節：一つの値の更新を行列で書く}

本編の例題$a_{n+1}=2a_n+1$では,定数$1$を別の成分に持たせると,
\[
  \begin{pmatrix}2&1\\0&1\end{pmatrix}
  \begin{pmatrix}a_n\\1\end{pmatrix}
  =\begin{pmatrix}2a_n+1\\1\end{pmatrix}
\]
と書ける.下の成分は毎回$1$のままで,上の成分だけが元の規則で更新される.

一般の$a_{n+1}=pa_n+q$ $(p\ne0,1)$でも,定数$q$を同じ更新に含めよう.\footnote{$p=0$なら,$n\ge2$で$a_n=q$と直接分かる.$p=1$は後で扱う.行列による更新の表現自体は,どちらの場合も使える.}
状態を$\boldsymbol{x}_n=\begin{pmatrix}a_n\\1\end{pmatrix}$とすると,
\[
  \begin{aligned}
    A&=\begin{pmatrix}p&q\\0&1\end{pmatrix},\qquad
    \boldsymbol{x}_{n+1}=A\boldsymbol{x}_n,\\
    \boldsymbol{x}_n&=A^{n-1}\boldsymbol{x}_1.
  \end{aligned}
\]
元の\bookterm{recurrence}{漸化式}と,変わらない成分$1$をまとめたものである.

本編で$b_n=a_n-\alpha$と置いた操作は,この状態から
\[
  b_n=(1,-\alpha)\boldsymbol{x}_n
\]
という量を取り出すことに当たる.
実際,\,$\alpha=p\alpha+q$を使えば
\[
  \begin{aligned}
    (1,-\alpha)A
      &=(p,q-\alpha)\\
      &=(p,-p\alpha)=p(1,-\alpha).
  \end{aligned}
\]
本編の$b_{n+1}=pb_n$は,\bookterm{left-eigen}{左固有ベクトル}で取り出した量の更新だった.
状態の分解として見ると,$\alpha=p\alpha+q$から
\[
  A\begin{pmatrix}\alpha\\1\end{pmatrix}
  =
  \begin{pmatrix}\alpha\\1\end{pmatrix}.
\]
\bookterm{fixed}{不動点}は更新で変わらない方向になる.一方,
\[
  A\begin{pmatrix}1\\0\end{pmatrix}
  =
  p\begin{pmatrix}1\\0\end{pmatrix}.
\]
\bookterm{fixed}{不動点}からのずれは,一歩ごとに $p$ 倍される方向にある.
実際,初期状態は
\[
  \begin{pmatrix}a_1\\1\end{pmatrix}
  =
  \begin{pmatrix}\alpha\\1\end{pmatrix}
  +(a_1-\alpha)\begin{pmatrix}1\\0\end{pmatrix}
\]
と分けられる.$n-1$回更新して上の成分を読めば,
\[
  a_n=\alpha+(a_1-\alpha)p^{n-1}.
\]
本編の「\bookterm{fixed}{不動点}」と「ずれ」が,倍率の異なる二方向になった.

この\bookterm{eigenvalue}{固有値}$1,p$は
\BookMathLayout{\[
  \det(\lambda I-A)
  =
  \det\begin{pmatrix}\lambda-p&-q\\0&\lambda-1\end{pmatrix}
  =(\lambda-p)(\lambda-1)=0
\]}{
\[
\begin{aligned}
\det(\lambda I-A)&=\det\begin{pmatrix}\lambda-p&-q\\0&\lambda-1\end{pmatrix}\\
&=(\lambda-p)(\lambda-1)=0
\end{aligned}
\]
}
からも得られる.本編の$\alpha=p\alpha+q$は不動点の\textbf{値},この\bookterm{characteristic}{特性方程式}は更新の\textbf{倍率}を求める式である.

本編では $p=1$ の場合を等差型として分けて扱った.
\bookterm{matrix}{行列}で書けば
\[
  A=\begin{pmatrix}1&q\\0&1\end{pmatrix},
  \qquad
  A^{n-1}=\begin{pmatrix}1&(n-1)q\\0&1\end{pmatrix}
\]
なので
\[
  a_n=a_1+(n-1)q
\]
となる.
$q\ne0$なら\bookterm{fixed}{不動点}はない.重なった\bookterm{eigenvalue}{固有値}$1$のもとで足し算が積み重なり,$n$が現れた.

\subsubsection{二回分の更新}\label{節：周期係数を行列で見直す}
  \sectionref*{節：周期的な係数}では,$2x+1$の後に$3x+2$を計算し,$6x+5$を得た.
  行列で書くと
  \[
    A=\begin{pmatrix}2&1\\0&1\end{pmatrix},\qquad
    B=\begin{pmatrix}3&2\\0&1\end{pmatrix}.
  \]
  先に$A$,次に$B$なので,二回分は
  \[
    BA=\begin{pmatrix}6&5\\0&1\end{pmatrix},\qquad
    BA\binom{x}{1}=\binom{6x+5}{1}
  \]
  となる.この例では$AB=BA$だが,一般には異なる.接続演習問4なら
  \[
    U=\begin{pmatrix}-1&0\\0&1\end{pmatrix},\qquad
    V=\begin{pmatrix}2&1\\0&1\end{pmatrix}
  \]
  に対して,
  \[
    VU=\begin{pmatrix}-2&1\\0&1\end{pmatrix},\qquad
    UV=\begin{pmatrix}-2&-1\\0&1\end{pmatrix}
  \]
  となる.更新を行う順序が,付加項の符号に現れた.
  一般の$a_{n+1}=p_na_n+q_n$でも,$\binom{a_{n+1}}1=
  \begin{pmatrix}p_n&q_n\\0&1\end{pmatrix}\binom{a_n}1$と表せる.
  一周期分は,先に行う\bookterm{matrix}{行列}を右に置いて掛け合わせる.途中の値は一回ずつの更新で求める.
  \bookproblemref{practice-connections}{4}{接続演習問4}で本文との対応を確かめよう.


\subsubsection{分数を二つの成分で追う}\label{節：一次分数型を行列で見直す}

\sectionref{節：1次分数型}で作った\bookterm{fixed}{不動点}からの差の比を,\bookterm{matrix}{行列}で見直す.
まず本編の
\[
  a_{n+1}=\frac{2a_n}{1-a_n}
\]
を使う.各項で$a_n\ne1$とし,\,$a_1\ne-1$とする.
一つの値を分子と分母に分けて持ち,$a_n=\dfrac{x_n}{y_n}$と書いてみよう.
$(x_1,y_1)=(a_1,1)$から
\[
  \begin{pmatrix}x_{n+1}\\y_{n+1}\end{pmatrix}
  =M\begin{pmatrix}x_n\\y_n\end{pmatrix},
  \qquad
  M=\begin{pmatrix}2&0\\-1&1\end{pmatrix}
\]
で更新すると,
\[
  \frac{x_{n+1}}{y_{n+1}}
  =\frac{2x_n}{y_n-x_n}
  =\frac{2a_n}{1-a_n}
\]
となり,元の更新に戻る.
本編で用いた$\displaystyle b_n=\frac{a_n}{a_n+1}$は,この二成分では
\[
  b_n=\frac{x_n}{x_n+y_n}
\]
である.分子は一回ごとに$2$倍されるが,分母は
\[
  x_{n+1}+y_{n+1}=2x_n+(y_n-x_n)=x_n+y_n
\]
と変わらない.従って,$b_{n+1}=2b_n$となる.
$x_1+y_1=a_1+1\ne0$なので,この比の分母はずっと零でない.

\bookterm{fixed}{不動点}$0,-1$を二成分で表すと,
\[
  M\begin{pmatrix}0\\1\end{pmatrix}
  =\begin{pmatrix}0\\1\end{pmatrix},
  \qquad
  M\begin{pmatrix}-1\\1\end{pmatrix}
  =2\begin{pmatrix}-1\\1\end{pmatrix}.
\]
対応する\bookterm{eigenvalue}{固有値}はそれぞれ$1,2$である.
差の比は,二方向の倍率の比を追っていた.

\bookstep{一般の一次分数型へ}
一般の係数でも,二成分の\bookterm{column-vector}{列ベクトル}で表せる.
\[
  a_{n+1}=\frac{pa_n+q}{ra_n+s}
\]
に対し,\,$r\ne0$, $ps-qr\ne0$を仮定し,各項で$ra_n+s\ne0$とする.
まず$(x_1,y_1)=(a_1,1)$とし,二つの数を
\[
  \begin{pmatrix}x_{n+1}\\y_{n+1}\end{pmatrix}
  =
  M\begin{pmatrix}x_n\\y_n\end{pmatrix},
  \qquad
  M=\begin{pmatrix}p&q\\r&s\end{pmatrix}
\]
で更新する.
$\displaystyle a_n=\frac{x_n}{y_n}$とすると
\[
  \frac{x_{n+1}}{y_{n+1}}
  =\frac{px_n+qy_n}{rx_n+sy_n}
  =\frac{pa_n+q}{ra_n+s}.
\]
また,\,$y_{n+1}=(ra_n+s)y_n$なので,もとの\bookterm{recurrence}{漸化式}の分母が零にならない限り,\,$y_n\ne0$が保たれる.
同じ零でない数で二成分を倍しても比は変わらない.この成分比を通して\bookterm{matrix}{行列}の作用を見る考え方を,射影変換という.

\bookterm{fixed}{不動点} $\alpha$ では
\[
  M\begin{pmatrix}\alpha\\1\end{pmatrix}
  =(r\alpha+s)\begin{pmatrix}\alpha\\1\end{pmatrix}.
\]
つまり\bookterm{fixed}{不動点}は\bookterm{eigenvalue}{固有値}ではなく,\bookterm{eigenvector}{固有ベクトル}の成分比である.
異なる二つの\bookterm{fixed}{不動点} $\alpha,\beta$ があり,対応する倍率を
$\lambda_\alpha=r\alpha+s,\ \lambda_\beta=r\beta+s$ とする.
状態を
\[
  \begin{pmatrix}x_n\\y_n\end{pmatrix}
  =c_n\begin{pmatrix}\alpha\\1\end{pmatrix}
  +d_n\begin{pmatrix}\beta\\1\end{pmatrix}
\]
と分けると,\,$\,c_n,d_n$ はそれぞれ $\lambda_\alpha,\lambda_\beta$ 倍される.
以下の比を取る際は$a_1\ne\beta$とする.\bookterm{fixed}{不動点}$\beta$を初項とする場合は$a_n=\beta$という定\bookterm{sequence}{数列}になる.
\bookterm{matrix}{行列}$M$は可逆で,\,\,$\lambda_\alpha\ne0$だから,\,\,$a_1\ne\beta$なら$c_n\ne0$,すなわち$a_n\ne\beta$が保たれる.
成分を比較すると
\[
  \begin{aligned}
    x_n&=\alpha c_n+\beta d_n,\qquad y_n=c_n+d_n,\\
    x_n-\alpha y_n&=(\beta-\alpha)d_n,\\
    x_n-\beta y_n&=(\alpha-\beta)c_n.
  \end{aligned}
\]
したがって,\,$\displaystyle a_n=\frac{x_n}{y_n}$より
\[
  \frac{a_n-\alpha}{a_n-\beta}=-\frac{d_n}{c_n},
  \qquad
  \frac{a_{n+1}-\alpha}{a_{n+1}-\beta}
  =\frac{\lambda_\beta}{\lambda_\alpha}
   \frac{a_n-\alpha}{a_n-\beta}.
\]
本編の\bookterm{fixed}{不動点}を求める方程式は$rx^2+(s-p)x-q=0$である.
解と係数の関係より$r(\alpha+\beta)=p-s$なので
\[
  \begin{aligned}
    \lambda_\beta&=r\beta+s=p-r\alpha,\\
    \lambda_\alpha&=r\alpha+s=p-r\beta.
  \end{aligned}
\]
よって公比は本編と同じ$\displaystyle \frac{\lambda_\beta}{\lambda_\alpha}=\frac{p-r\alpha}{p-r\beta}$になる.
置き換えは二つの固有方向の成分比を取り出す操作であり,最初の例は倍率$1,2$の場合だった.

\subsection{二つの値を更新する}\label{節：状態を二つ持つと,二階の漸化式になる}

次は,二つの数を同時に更新する規則を考える.例えば
\[
  x_{n+1}=x_n+y_n,
  \qquad
  y_{n+1}=x_n
\]
という二つの規則は
\[
  \begin{pmatrix}x_{n+1}\\y_{n+1}\end{pmatrix}
  =
  \begin{pmatrix}1&1\\1&0\end{pmatrix}
  \begin{pmatrix}x_n\\y_n\end{pmatrix}
\]
と一つの\bookterm{matrix}{行列}で表せる.
初めを $(x_0,y_0)=(1,0)$ とすると,
\[
  (1,0),\ (1,1),\ (2,1),\ (3,2),\ldots
\]
と進み,\,$\,x_n=F_{n+1},\,y_n=F_n$ となる.
ここで $F_0=0,\,F_1=1$ はフィボナッチ\bookterm{sequence}{数列}である.
一歩前の$x$を$y$に残すので,$x$だけを見れば
\[
  x_{n+2}=x_{n+1}+x_n
\]
という二階の\bookterm{recurrence}{漸化式}になる.

\sectionref{節：隣接3項間漸化式型}も,直前の二項を状態にすればよい.一般に
\[
  a_{n+2}=s a_{n+1}+t a_n
\]
に対して
\[
  \boldsymbol{v}_n=
  \begin{pmatrix}a_{n+1}\\a_n\end{pmatrix},
  \qquad
  B=\begin{pmatrix}s&t\\1&0\end{pmatrix}
\]
とすると,
\[
  B\boldsymbol{v}_n
  =
  \begin{pmatrix}s a_{n+1}+t a_n\\a_{n+1}\end{pmatrix}
  =
  \begin{pmatrix}a_{n+2}\\a_{n+1}\end{pmatrix}
  =\boldsymbol{v}_{n+1}.
\]
二行目が直前の項を残すので,$\boldsymbol{v}_{n+1}=B\boldsymbol{v}_n$という一階更新になる.
なお本編の形 $a_{n+2}+pa_{n+1}+qa_n=0$ で書く場合は,ここで $s=-p,\,t=-q$ とすればよい.

$B\begin{pmatrix}x\\y\end{pmatrix}=\lambda\begin{pmatrix}x\\y\end{pmatrix}$の二行目から$x=\lambda y$となる.
零でない\bookterm{eigenvector}{固有ベクトル}では$y\ne0$なので,$y=1$にそろえると$\begin{pmatrix}\lambda\\1\end{pmatrix}$になる.一行目から
\[
  s\lambda+t=\lambda^2,
  \qquad
  \lambda^2-s\lambda-t=0
\]
を得る.
これは$\lambda\ne0$なら,$a_n=\lambda^n$をもとの\bookterm{recurrence}{漸化式}に代入し,共通の冪を除いて得る\bookterm{characteristic}{特性方程式}と一致する.
$\lambda=0$も\bookterm{matrix}{行列}ではそのまま扱える.\sectionref*{節：行列の基本}で見たように,零の固有方向は初期状態にだけ寄与し,一回の更新で消える.
本編の符号で書けば $\lambda^2+p\lambda+q=0$ である.

この条件は\bookterm{linear}{線形}代数では
\BookMathLayout{\[
  \det(\lambda I-B)
  =
  \det\begin{pmatrix}\lambda-s&-t\\-1&\lambda\end{pmatrix}
  =\lambda^2-s\lambda-t=0
\]}{
\[
\begin{aligned}
\det(\lambda I-B)&=\det\begin{pmatrix}\lambda-s&-t\\-1&\lambda\end{pmatrix}\\
&=\lambda^2-s\lambda-t=0
\end{aligned}
\]
}
とも書く.
本編の\bookterm{root}{特性根}は,更新行列の\bookterm{eigenvalue}{固有値}だった.

二つの異なる\bookterm{root}{特性根} $\lambda_1,\lambda_2$ があるとき,対応する\bookterm{eigenvector}{固有ベクトル}
\[
  \begin{pmatrix}\lambda_1\\1\end{pmatrix},
  \qquad
  \begin{pmatrix}\lambda_2\\1\end{pmatrix}
\]
は同じ直線上にはない.
初めの状態を二つの方向に分けておけば,\,$k$回の更新でそれぞれの部分が $\lambda_1^k,\,\lambda_2^k$ 倍される.
こうして\bookterm{general}{一般項}に\bookterm{root}{特性根}の冪が現れる.

さらに,本編で作った$a_{n+1}-\lambda_1a_n$という量も,同じ\bookterm{matrix}{行列}の中に見つかる.
解と係数の関係$\lambda_1+\lambda_2=s$, $\lambda_1\lambda_2=-t$を使うと
\[
  \begin{aligned}
    (1,-\lambda_1)B
      &=(s-\lambda_1,t)\\
      &=(\lambda_2,-\lambda_1\lambda_2)
       =\lambda_2(1,-\lambda_1).
  \end{aligned}
\]
したがって
\[
  \begin{aligned}
    b_n&=(1,-\lambda_1)\boldsymbol{v}_n
         =a_{n+1}-\lambda_1a_n,\\
    b_{n+1}&=\lambda_2 b_n.
  \end{aligned}
\]
根を入れ替えればもう一方も得られる.本編の$a_{n+2}=5a_{n+1}-6a_n$で作った$a_{n+1}-2a_n$は,$(1,-2)$で取り出す$3$倍の量だった.
\bookterm{eigenvector}{固有ベクトル}で状態を分ける方法と,\bookterm{left-eigen}{左固有ベクトル}で量を取り出す方法は,同じ更新の二通りの見方である.
二つの根が複素数でも,成分を複素数まで広げて同じ計算を行える.
実数の初期値から出発した場合,複素\bookterm{conjugate}{共役}な二方向の寄与が組になり,\bookterm{sequence}{数列}の項は実数のままになる.

\sectionref{節：連立型}の$A=\begin{pmatrix}2&1\\1&2\end{pmatrix}$でも,和と差$u_n=x_n+y_n$, $v_n=x_n-y_n$は
\[
  u_{n+1}=3u_n,\qquad v_{n+1}=v_n
\]
を満たす.前節の左固有ベクトル$(1,1)$は和を,$(1,-1)$は差を取り出していた.

図の$(1,1)$と$(1,-1)$の方向を使うと,座標軸を$45^\circ$回して見ることになる.
回すのは\textbf{座標の見方}であり,更新行列$A$自体が回転行列なのではない.
  \BookFigSystem

向きをそろえて回転\bookterm{matrix}{行列}
\[
  Q=\frac{1}{\sqrt2}
  \begin{pmatrix}
    1&-1\\
    1&1
  \end{pmatrix}
\]
を用いると,
\[
  Q^{\mathsf T}
  \begin{pmatrix}x_n\\y_n\end{pmatrix}
  =\frac{1}{\sqrt2}
  \begin{pmatrix}u_n\\-v_n\end{pmatrix}
\]
となる.
$\boldsymbol z_n=\begin{pmatrix}x_n\\y_n\end{pmatrix}$とし,新しい座標を$\boldsymbol w_n=Q^{\mathsf T}\boldsymbol z_n$と書く.
直接掛けると$Q^{\mathsf T}Q=QQ^{\mathsf T}=I$だから,元へ戻す式は$\boldsymbol z_n=Q\boldsymbol w_n$である.
元の更新$\boldsymbol z_{n+1}=A\boldsymbol z_n$にこれを入れると,
\[
  \boldsymbol w_{n+1}
  =Q^{\mathsf T}\boldsymbol z_{n+1}
  =Q^{\mathsf T}AQ\boldsymbol w_n.
\]
元の座標へ戻す$Q$,更新する$A$,新座標へ直す$Q^{\mathsf T}$の順に作用している.
ここで
\[
  Q^{\mathsf T}AQ=
  \begin{pmatrix}3&0\\0&1\end{pmatrix}
\]
と表される.
だから,新座標では二つの等比型に分かれる.

実際に計算へ戻せば,
\[
  \begin{aligned}
    u_n&=3^{n-1}(x_1+y_1),\\
    v_n&=x_1-y_1,\\
    x_n&=\frac{u_n+v_n}{2},\qquad
    y_n=\frac{u_n-v_n}{2}
  \end{aligned}
\]
となる.
$x_1=1$, $y_1=0$では,前節の$A^k\begin{pmatrix}1\\0\end{pmatrix}$と$k=n-1$で一致する.
斜めの二方向を更新し,元の座標へ戻した結果である.

一般に連立\bookterm{recurrence}{漸化式}を
\BookMathLayout{\[
  \boldsymbol{z}_{n+1}=A\boldsymbol{z}_n,
  \qquad
  \boldsymbol{z}_n=\begin{pmatrix}a_n\\b_n\end{pmatrix},
  \qquad
  A=\begin{pmatrix}p&q\\r&s\end{pmatrix}
\]}{
\[
\begin{aligned}
\boldsymbol z_{n+1}&=A\boldsymbol z_n,\\
\boldsymbol z_n&=\begin{pmatrix}a_n\\b_n\end{pmatrix},\quad
A=\begin{pmatrix}p&q\\r&s\end{pmatrix}
\end{aligned}
\]
}
とし,零でない\bookterm{column-vector}{列ベクトル}$\boldsymbol{\ell}$を用いて$\,c_n=\boldsymbol{\ell}^{\mathsf T}\boldsymbol{z}_n$ とおく.
もし
\[
  \boldsymbol{\ell}^{\mathsf T}A
  =\gamma\boldsymbol{\ell}^{\mathsf T}
\]
なら,\,$\,c_{n+1}=\gamma c_n$ となる.
本編の連立型で $c_n=\alpha a_n+\beta b_n$ が\bookterm{geometric}{等比数列}となるよう $\alpha,\beta$ を選ぶ計算は,\bookterm{left-eigen}{左固有ベクトル}を求めていたのである.
この関係を満たす零でないベクトル$\boldsymbol{\ell}$が存在するための条件
\[
  \det(\gamma I-A)
  =\gamma^2-(p+s)\gamma+(ps-qr)=0
\]
は,本編で現れた\bookterm{characteristic}{特性方程式}と一致する.
異なる二つの\bookterm{eigenvalue}{固有値}があれば,対応する\bookterm{left-eigen}{左固有ベクトル}は比例しないので,取り出した二つの量から$a_n,b_n$を求め直せる.
二階の\bookterm{recurrence}{漸化式}も二成分の状態に直せば,この枠組みに入る.
フィボナッチ\bookterm{sequence}{数列}では $s=t=1$ なので
\[
  B=\begin{pmatrix}1&1\\1&0\end{pmatrix}
\]
であり,\bookterm{root}{特性根}は
\[
  \varphi=\frac{1+\sqrt5}{2},
  \qquad
  \psi=\frac{1-\sqrt5}{2}
\]
である.
初期値 $F_0=0,\,F_1=1$ に合わせて二つの方向の割合を決めると
\[
  F_n=\frac{\varphi^n-\psi^n}{\sqrt5}
\]
を得る.
$|\psi|<1<\varphi$だから,$\psi^n$は符号を交互に変えながら減衰し,$\varphi^n$が増大を支配する.
\bookterm{root}{特性根}は長期的な動きも表す.

\bookterm{characteristic}{特性方程式}が重解 $r\ne0$ を持つ場合には,二つの異なる倍率の方向に分けられない.
\bookterm{general}{一般項}に $n$ が現れる理由は,\bookterm{recurrence}{漸化式}に直接戻ると見やすい.
重解のときは
\[
  a_{n+2}=2r a_{n+1}-r^2a_n
\]
だから,\,$\,a_n=r^n b_n$ と置いて代入すると
\[
  b_{n+2}-2b_{n+1}+b_n=0
\]
となる.
差が一定だから$b_n=C+Dn$,つまり$a_n=(C+Dn)r^n$となる.
\sectionref{節：二次特性方程式型}の重解の例題で得た$a_n=n2^{n-1}$も,等比的な倍率$2^{n-1}$を取り除くと$n$という等差的な増加が残る例である.
\sectionref{節：階差等比複合型}の指数型で倍率が一致した場合とも,等比部分を取り除いて差を取る工夫が共通する.
たとえば重解が $1$ のときは
\BookMathLayout{\[
  a_{n+2}-2a_{n+1}+a_n
  =(a_{n+2}-a_{n+1})-(a_{n+1}-a_n)=0
\]}{
\[
\begin{aligned}
a_{n+2}-2a_{n+1}+a_n\\
{}=(a_{n+2}-a_{n+1})-(a_{n+1}-a_n)\\
{}=0
\end{aligned}
\]
}
より,隣り合う項の差が一定で $a_n=C+Dn$ となる.
重解が $r=0$ の場合は $a_{n+2}=0$ となる退化した\bookterm{recurrence}{漸化式}であり,直接調べればよい.

\subsection{角度の整数倍と多項式}\label{節：三角関数型漸化式からチェビシェフ多項式へ}

\sectionref{節：虚数解型隣接3項間漸化式}では,複素数の冪と正弦・余弦という\bookterm{basis}{基底}を使い分けた.
\sectionref{節：三角関数型}では,値の更新を角度の更新に直した.
この二つを「整数倍の角を表す多項式」でつなぐ.

まず,一定の角$\theta$だけ回す操作を考える.
\bookterm{addition}{加法定理}より
\[
  R=\begin{pmatrix}\cos\theta&-\sin\theta\\
                  \sin\theta&\cos\theta\end{pmatrix}
\]
とおけば
\[
  \begin{pmatrix}\cos((n+1)\theta)\\\sin((n+1)\theta)\end{pmatrix}
  =R\begin{pmatrix}\cos(n\theta)\\\sin(n\theta)\end{pmatrix}.
\]
この\bookterm{matrix}{行列}の\bookterm{characteristic}{特性方程式}は
\[
  \lambda^2-2\cos\theta\,\lambda+1=0
\]
であり,\,$0<\theta<\pi$なら\bookterm{eigenvalue}{固有値}は$\cos\theta\pm i\sin\theta$である.
複素数の冪による回転が,\bookterm{eigenvalue}{固有値}にも現れた.

横座標だけを追うと,余弦の\bookterm{addition}{加法定理}より
\[
  a_{n+2}=2\cos\theta\,a_{n+1}-a_n
\]
という隣接3項間\bookterm{recurrence}{漸化式}になる.
ここでは\textbf{一定角を足す}.本編の\bookterm{trigonometric}{三角関数}型の\textbf{角度を二倍にする}更新とは異なるが,以下の多項式はどちらの整数倍角も表せる.

$x=\cos\theta$とすると,\bookterm{addition}{加法定理}から
\[
  \cos((n+1)\theta)=2x\cos(n\theta)-\cos((n-1)\theta).
\]
同じ規則で多項式の列を
\begin{align*}
  &T_0(x)=1,\qquad T_1(x)=x,\\
  &T_{n+1}(x)=2xT_n(x)-T_{n-1}(x)
  \qquad(n\ge1)
\end{align*}

と定める.
最初から足し算と掛け算だけで作るので,すべて多項式になる.これを\bookterm{chebyshev}{チェビシェフ多項式}という.

この定義から
\[
  T_2(x)=2x^2-1,\qquad
  T_3(x)=4x^3-3x
\]
が得られる.
二倍角・三倍角の公式と同じ形である.一般にも,$n=0,1$の値と更新規則が一致するので,帰納法から
\[
  T_n(\cos\theta)=\cos(n\theta).
\]
整数倍の角の余弦を,$x=\cos\theta$の多項式で表せた.図は$T_2$による二倍角の横座標を示す.
  \BookFigChebyshev

ここで本編の$\displaystyle a_1=\frac{\sqrt3}{2}$, $a_{n+1}=2a_n^2-1$に戻ろう.
本編では$\displaystyle a_1=\cos(\frac{\pi}{6})$から角度を二倍ずつ進めたが,多項式の言葉では
\[
  a_{n+1}=T_2(a_n),\qquad
  a_n=T_{2^{n-1}}(a_1)
\]
と書ける.
初項そのものの多項式で\bookterm{general}{一般項}を表せた.その理由は次の合成の関係にある.

一般に,$x=\cos\theta$なら
\BookMathLayout{\[
  T_m(T_k(x))
  =T_m(\cos(k\theta))
  =\cos(mk\theta)
  =T_{mk}(x).
\]}{
\[
\begin{aligned}
T_m(T_k(x))&=T_m(\cos(k\theta))\\
&=\cos(mk\theta)=T_{mk}(x).
\end{aligned}
\]
}
この等式は $[-1,1]$ のすべての $x$ で成り立ち,両辺が多項式なので,多項式の恒等式として実数全体で成り立つ.
したがって$T_2$を繰り返し合成する操作は,添字を二倍ずつ進める$T_2,T_4,T_8,\ldots$にまとまる.
本編の「整数倍の公式へ」で扱った$m$倍角にも同じことがいえる.
整数$m\ge2$に対して$a_{n+1}=T_m(a_n)$なら
\[
  a_n=T_{m^{n-1}}(a_1)
\]
である.
$-1\le a_1\le1$なら$a_1=\cos\theta$と表せるため,これは本編の$a_n=\cos(m^{n-1}\theta)$と一致する.
多項式としての式は,\bookterm{trigonometric}{三角関数}で表せない初項についても成り立つ.

多項式列を\bookterm{recurrence}{漸化式}で作る例は,\bookterm{chebyshev}{チェビシェフ多項式}だけではない.
\bookterm{legendre}{ルジャンドル多項式} $P_n(x)$ は,初期値
\[
  P_0(x)=1,\qquad P_1(x)=x
\]
と三項\bookterm{recurrence}{漸化式}
\BookMathLayout{
\[
\begin{aligned}
(n+1)P_{n+1}(x)&=(2n+1)xP_n(x)\\
&\quad-nP_{n-1}(x)\qquad(n\ge1)
\end{aligned}
\]
}{
\[
\begin{aligned}
(n+1)P_{n+1}(x)&=(2n+1)xP_n(x)\\
&\quad-nP_{n-1}(x)\qquad(n\ge1)
\end{aligned}
\]
}
で定まる.
例えば
\[
  P_2(x)=\frac{3x^2-1}{2}
\]
である.
\bookterm{chebyshev}{チェビシェフ多項式}と同じく三項\bookterm{recurrence}{漸化式}から作られるが,\bookterm{legendre}{ルジャンドル多項式}では係数が $n$ にも依存する.

\bookterm{legendre}{ルジャンドル多項式}は,区間 $[-1,1]$ 上で重み $1$ に関して直交する.
ここでいう直交性とは,異なる次数の二つを掛けて区間全体で積分すると $0$ になる性質である.
\[
  \int_{-1}^{1}P_m(x)P_n(x)\,dx=0
  \qquad(m\ne n)
\]
この性質は関数の表示や近似に役立つ.ここでは,三項\bookterm{recurrence}{漸化式}が生む別の多項式列として紹介する.

\subsection*{コラム：有限個の点から多項式を作る}

最初の三項が $1,2,3$ なら,
$x$ と $x+(x-1)(x-2)(x-3)$ はどちらもその三項に一致するが,
$n=4$ ではそれぞれ $4$ と $10$ になる.
有限個の値が一致しても,その後の値まで同じとは限らない.
与えられた値を通る多項式を,一般の場合にも作ってみよう.

最初の$N$項$a_1,\ldots,a_N$に対して,
\[
  L(x)=\sum_{k=1}^{N}a_k
  \prod_{\substack{1\le j\le N\\j\ne k}}
  \frac{x-j}{k-j}
\]
と置くと, $L(k)=a_k$ $(1\le k\le N)$ となる.
$x=k$ のとき $k$ 番目の積は $1$ になり,ほかの積には $x-k$ の因子があるためである.
次数を $N-1$ 以下に制限した多項式の中では,この $L$ は一意である.二つの候補の差は $1,\ldots,N$ を根にもつ次数 $N-1$ 以下の多項式なので,零多項式になる.

次数制限はデータからは決まらない.$L(x)+c\prod_{j=1}^{N}(x-j)$も,任意の定数$c$で最初の$N$個の値を保つ.
次の図は,別の三点を通り,途中の形が異なる二つの曲線を描いた例である.
  \BookFigLagrange

\bookterm{interpolation}{ラグランジュ補間}は一つの延長候補を作る.それを元の規則と結論づけるには,次数や\bookterm{recurrence}{漸化式}など追加の条件が要る.

\subsection{不動点付近の線形化}\label{節：非線形な更新も,不動点の近くでは線形に見える}

\sectionref*{節：特性方程式型}では,不動点からのずれが$p$倍され,$|p|<1$なら近づいた.
\bookterm{nonlinear}{非線形}の更新でも,微分係数と一次近似で不動点付近の倍率を読める.
一項前から次が決まる\bookterm{recurrence}{漸化式}
\[
  a_{n+1}=f(a_n)
\]
を考え,\,$\,\alpha$ を\bookterm{fixed}{不動点} $f(\alpha)=\alpha$ とする.
そこからのずれを $e_n=a_n-\alpha$ と置くと
\[
  e_{n+1}=f(\alpha+e_n)-f(\alpha)
\]
である.
$f$ が $\alpha$ で微分可能なら,微分係数の定義から
\[
  \frac{f(\alpha+e)-f(\alpha)}{e}\longrightarrow f'(\alpha)
  \qquad(e\longrightarrow0)
\]
である.
つまり,ずれ $e$ が十分小さい範囲では
\[
  f(\alpha+e)-f(\alpha)\approx f'(\alpha)e
\]
となる.
不動点の近くで曲線を接線に置き換えたのである.収束も調べられる.
もし $|f'(\alpha)|<1$ なら,\,$\,|f'(\alpha)|$ より大きく $1$ より小さい数 $c$ を選べる.
微分係数の定義により,\,$\,a_n$ が $\alpha$ に十分近ければ
\[
  |e_{n+1}|
  =|f(\alpha+e_n)-f(\alpha)|
  \le c|e_n|
\]
となる.
誤差が毎回$c$倍以下に縮むので,十分近い初期値は近くにとどまり,$\alpha$へ収束する.
行列の\bookterm{eigenvalue}{固有値}と同じく,$f'(\alpha)$が局所的なずれの倍率を表す.
一次分数型で先ほどの二つの異なる\bookterm{fixed}{不動点}がある場合には
\[
  \begin{aligned}
    f'(\alpha)
      &=\frac{ps-qr}{(r\alpha+s)^2}\\
      &=\frac{\lambda_\alpha\lambda_\beta}{\lambda_\alpha^2}
       =\frac{\lambda_\beta}{\lambda_\alpha}.
  \end{aligned}
\]
特性方程式から$ps-qr=\lambda_\alpha\lambda_\beta$である.
比の置き換えで得た\textbf{正確な公比}と,微分による\textbf{局所的な倍率}が一致した.
\bookterm{linear}{線形}代数と微分は,一歩の変化を倍率で測る見方でつながる.

\subsection{整数から作る近似分数}\label{節：ペル方程式と近似分数}
  \sectionref*{節：ニュートン法型}では,$a_1=2$から$\sqrt2$を近似した.
  ここでは同じ更新式$a_{n+1}=\frac12\left(a_n+\dfrac{2}{a_n}\right)$で,初項を$a_1=1$に変えてみる.
  \[
    a_2=\frac32,\qquad a_3=\frac{17}{12},\qquad a_4=\frac{577}{408}
  \]
  となる.初項$1,2$のどちらも$\dfrac{3}{2}$へ移るため,第2項以降は本編と一致する.
  分子と分母を追い,連立型・一次分数型・\bookterm{newton}{ニュートン法}をつなげる.微積分は使わない.

  \bookstep{分子と分母の関係}
  $3^2-2\cdot2^2=17^2-2\cdot12^2=577^2-2\cdot408^2=1$である.
  次の\bookterm{sequence}{数列}を作ろう.
  \[
    x_k+y_k\sqrt2=(3+2\sqrt2)^k\quad(k\ge1)
  \]
  $3+2\sqrt2$倍して係数を比べると,
  \[
    x_{k+1}=3x_k+4y_k,\qquad y_{k+1}=2x_k+3y_k
  \]
  となる.$x_1=3$, $y_1=2$なので,両方とも正の整数である.
  \bookterm{conjugate}{共役}な式との積から
  \[
    x_k^2-2y_k^2=(3+2\sqrt2)^k(3-2\sqrt2)^k=1
  \]
  が保たれる.整数の式を保つ更新は,
  \bookproblemref{practice-bridge}{9}{発展入門問9}の見方とも共通する.

  \bookstep{比から一次分数型へ}
  $r_k=\frac{x_k}{y_k}$とすると,連立の更新から
  \[
    r_{k+1}=\frac{3r_k+4}{2r_k+3}
  \]
  である.分母は正であり,\bookterm{fixed}{不動点}を求めると$\pm\sqrt2$になる.
  一次分数型で用いた差の比は,
  \[
    \frac{r_{k+1}-\sqrt2}{r_{k+1}+\sqrt2}
    =(3-2\sqrt2)^2
      \frac{r_k-\sqrt2}{r_k+\sqrt2}
  \]
  となる.整数の組とその比は,同じ更新の二つの表現である.

  \bookstep{例：$\dfrac{17}{12}$の次の近似分数}
  本編で使った更新を$\dfrac{17}{12}$に施すと,
  \[
    \frac12\left(\frac{17}{12}+\frac{24}{17}\right)=\frac{577}{408}
  \]
  となる.一方,整数の更新を一回進めた$r_3$は$\dfrac{99}{70}$である.
  ニュートン法は,この列で次の分数へ一つ進んでいるわけではない.

  \bookstep{別解：ニュートン法と添字の倍増}
  $(x_k+y_k\sqrt2)^2$の展開から,
  \[
    x_{2k}=x_k^2+2y_k^2,\qquad y_{2k}=2x_ky_k
  \]
  なので,
  \[
    r_{2k}=\frac12\left(r_k+\frac2{r_k}\right).
  \]
  この\bookterm{newton}{ニュートン法}型の式から
  \[
    a_n=r_{2^{n-2}}\qquad(n\ge2)
  \]
  となる.$\dfrac{3}{2}$, $\dfrac{17}{12}$, $\dfrac{577}{408}$は$r_1$, $r_2$, $r_4$であり,連続する三項ではない.
  \BookFigPellDoubling
  $x_k^2-2y_k^2=1$を因数分解すれば,
  \[
    r_k-\sqrt2=\frac1{y_k^2(r_k+\sqrt2)}>0.
  \]
  $y_{k+1}>3y_k$なので$y_k\to\infty$であり,$r_k\to\sqrt2$となる.
  添字の倍増からも,ニュートン法の速さを読める.

  \bookstep{ペル方程式}
  平方数でない正の整数$d$に対し,$x^2-dy^2=1$の整数解を求める問題を\textbf{\bookterm{pell}{ペル方程式}}という.\footnote{ここでは$d=2$の具体例を扱い,一般の解の存在や全解の分類には進まない.}

  \bookstep{連分数にも現れる分数}
  $\sqrt2=1+\frac1{1+\sqrt2}$から,有限段階の近似を
  \[
    C_0=1,\qquad C_{j+1}=1+\frac1{1+C_j}
    \quad(j\ge0)
  \]
  で定める.$C_1=\dfrac{3}{2}$, $C_2=\dfrac{7}{5}$, $C_3=\dfrac{17}{12}$である.
  $C_j=\dfrac{P_j}{Q_j}$とすると,$P_0=Q_0=1$から
  \[
    P_{j+1}=P_j+2Q_j,\qquad Q_{j+1}=P_j+Q_j
  \]
  が得られ,$P_j+Q_j\sqrt2=(1+\sqrt2)^{j+1}$となる.
  $(1+\sqrt2)^2=3+2\sqrt2$だから,
  \[
    C_{2k-1}=r_k\qquad(k\ge1)
  \]
  である.この有限段階の値を\textbf{\bookterm{continued}{連分数}の近似分数}と呼ぶ.
  偶数番号では$P_j^2-2Q_j^2=-1$,奇数番号では$1$となり,
  同じ誤差の因数分解で$C_j\to\sqrt2$も確認できる.

  \begin{bookpoint}
    \textbf{三つの解法のつながり}\par
    分子・分母を追えば連立型,比を追えば一次分数型,
    添字を倍に進めれば平方根の\bookterm{newton}{ニュートン法}になる.
    追う量と更新回数を変えると,解法がつながる.
  \end{bookpoint}
  保存される式と誤差を\bookproblemref{practice-connections}{6}{接続演習問6}で確かめよう.
  \bookterm{pell}{ペル方程式}と\bookterm{continued}{連分数}の補足資料は,\chapterref*{章：参考文献}のKeith Conradの二つの解説にある.

